\diarytitle{0079}{パラメータを含む3次正方行列が正則でない条件}

$A(t)$ が正則でないための必要十分条件は $\det A(t) = 0$ であるから、まず $\det A(t)$ を計算する。第 2 行・第 3 行からそれぞれ第 1 行を引いても行列式の値は変わらないから、
\[
\det A(t) =
\begin{vmatrix}
t & 1 & 1 \\
1-t & t-1 & 0 \\
1-t & 0 & t-1
\end{vmatrix}
\]
第 2 行・第 3 行から共通因子 $(t-1)$ をくくり出すと
\[
\det A(t) = (t-1)^{2}
\begin{vmatrix}
t & 1 & 1 \\
-1 & 1 & 0 \\
-1 & 0 & 1
\end{vmatrix}
\]
右辺の $3\times3$ 行列式を第 1 行で余因子展開すると
\[
\begin{vmatrix}
t & 1 & 1 \\
-1 & 1 & 0 \\
-1 & 0 & 1
\end{vmatrix}
= t(1\cdot1 - 0\cdot0) - 1\bigl((-1)\cdot1 - 0\cdot(-1)\bigr) + 1\bigl((-1)\cdot0 - 1\cdot(-1)\bigr)
= t + 1 + 1 = t+2
\]
したがって
\[
\det A(t) = (t-1)^{2}(t+2)
\]
（検算: $t=0$ を代入すると $\det A(0) = (-1)^2\cdot 2 = 2$。一方、直接展開すると $A(0)=\begin{pmatrix}0&1&1\\1&0&1\\1&1&0\end{pmatrix}$ で、$\det = 0(0-1)-1(0-1)+1(1-0)=0+1+1=2$ となり一致する。）

$A(t)$ が正則でないのは $\det A(t) = 0$ のときであり、$(t-1)^{2}(t+2)=0$ を解くと
\[
\boxed{\; t = 1 \ \text{（重解）}, \quad t = -2 \;}
\]
検算として、$t=1$ のとき $A(1)=\begin{pmatrix}1&1&1\\1&1&1\\1&1&1\end{pmatrix}$ は全行が等しく階数 $1$ で正則でないことが確かめられる。$t=-2$ のとき $A(-2)=\begin{pmatrix}-2&1&1\\1&-2&1\\1&1&-2\end{pmatrix}$ は各行の和が $0$ となるため $(1,1,1)^{T}$ が核に属し、正則でないことが確かめられる。
